\uuid{Wilv}
\exo7id{1703}
\titre{exo7 1703}
\auteur{barraud}
\organisation{exo7}
\datecreate{2003-09-01}
\isIndication{false}
\isCorrection{false}
\chapitre{Réduction d'endomorphisme, polynôme annulateur}
\sousChapitre{Applications}
\cpge{mpsi-mp2i.calcul-matriciel-systemes.e, pcsi.determinants.c, ptsi.determinants.a, mpsi-mp2i.polynomes-fractions-rationnelles.a, pcsi.polynomes.a, ptsi.polynomes.a, bcpst1.matrices, mp-mpi.reduction-endomorphismes.g, pc.complements-espaces-vectoriels.d, psi.complements-espaces-vectoriels.d, mp-mpi.reduction-endomorphismes.c, pc.reduction.b, psi.reduction-endomorphismes.b, pt.reduction-endomorphismes.b}
\module{Algèbre}
\niveau{L2}
\difficulte{}

\contenu{
\question{
Soit $A\in\mathcal{M}_{3}(\R)$ la matrice 
  $$
  A=
  \begin{pmatrix}
    0&1&1\\
    1&0&1\\
    0&0&1
  \end{pmatrix}
  $$
  Donner un polynôme annulateur de $A$ de degré aussi petit que
  possible. En déduire $A^{-1}$, $A^{3}$, et $A^{5}$.
}
}
